% sections/model.tex

\subsection{Model Architecture}

The model is a stochastic Petri net consisting of 28 places and 29 transitions, specified in the Signal Hierarchical Petri Nets model description language and provided in two model files: \texttt{phase3a\_spatial\_scaled.shy} (Phase~B; 0.01$\times$ initial conditions for all internal species) and \texttt{phase3a\_spatial\_clean.shy} (Phase~C; 1$\times$ standard initial conditions). The two files share identical topology, arc structure, and rate expressions; they differ only in the initial markings of the intracellular species. Places represent molecular species partitioned into six biological compartments: extracellular cytokines (EPO, GCSF), membrane receptors (EPOR, GCSFR, and their bound forms, internalized forms, and recycling pools), cytoplasmic and nuclear mRNA for both GATA1 and PU.1, cytoplasmic and nuclear protein for both factors, phosphorylated nuclear GATA1 (pGATA1\_nuc), and the energy currency species (ATP, ADP, GTP, GDP, Pi). Transitions represent the reactions that connect these species: receptor binding and unbinding, receptor internalisation and recycling, mRNA transcription, nuclear export, mRNA degradation, protein translation, nuclear import, protein degradation, GATA1 phosphorylation, and ubiquitin-mediated PU.1 degradation. Energy metabolism is represented through ATP synthesis (from ADP and Pi) and GTP regeneration (from GDP and Pi), ensuring that the energy currency species are dynamically coupled to the signalling network rather than held at fixed concentrations.

The complete Petri net topology is shown in Fig.~\ref{fig:model}.

\begin{figure}[!ht]
  \centering
  \includegraphics[width=\linewidth]{phase3a_spatial_scaled}
  \caption{\textbf{Stochastic Petri net topology of the GATA1/PU.1 lineage commitment model.} 28 places (species) and 29 transitions (reactions) spanning six biological compartments. Node labels show species names and initial concentrations at the Phase B single-cell scale (0.01$\times$ standard; see Section~\ref{sec:methods}); the arc topology and rate expressions are identical in \texttt{phase3a\_spatial\_clean.shy}.}
  \label{fig:model}
\end{figure}

The model's topology implements two opposing regulatory arms. In the erythroid arm, EPOR\_bound promotes GATA1\_transcription by appearing in the rate expression as an activating factor. The resulting GATA1 mRNA is exported to the cytoplasm, translated, imported into the nucleus, and phosphorylated (at the residue corresponding to Ser310 in the human protein) by a kinase whose activity scales with EPOR\_bound. The phosphorylated form, pGATA1\_nuc, then drives PU1\_ubiquitin\_degradation---a transition whose rate is proportional to the product of pGATA1\_nuc abundance and PU1\_Protein\_nuc abundance, representing the GATA1-dependent recruitment of an E3 ubiquitin ligase to PU.1. In the myeloid arm, GCSFR\_bound promotes PU1\_transcription by an analogous activating factor in its rate expression. PU.1 protein in the nucleus suppresses GATA1\_transcription through a competitive inhibition term that reduces the effective rate of that transition. These kinetic forms produce the double-negative feedback topology necessary for bistability.

\subsection{Kinetic Parameters and Calibration}

Rate constants were initially derived from published measurements for EPOR binding kinetics~\cite{huang2007}, GATA1 and PU.1 transcription and degradation rates~\cite{chickarmane2009}, and receptor internalisation and recycling dynamics. Energy metabolism parameters were calibrated to maintain physiologically reasonable ATP and GTP levels under steady-state conditions, with GTP regeneration and ATP synthesis rate constants scaled to produce ATP+ADP $\approx$ 33~mM and GTP+GDP $\approx$ 17.73~mM---values consistent with published measurements for erythroid progenitor cells. Iterative calibration confirmed that the energy system remained stable across the full range of simulated conditions and did not limit or bias commitment outcomes.

The cross-repression strengths and the basal transcription rates were calibrated to produce, at population-scale concentrations (1$\times$ initial conditions), a sharp commitment threshold near EPO = 0.449~mM---consistent with the EPO* observed in bulk colony assays for EPO concentrations in the range of 0.1--1~U/mL. At this scale the model exhibits deterministic-appearing dynamics with very low replicate variance. The critical test of the model was its behaviour at reduced copy numbers, which we address through the 0.01$\times$ scaling strategy described below.

\subsection{Single-Cell Copy Number Scaling}

To access the stochastic regime corresponding to physiological single-cell copy numbers, the initial markings of all internal molecular species (mRNA, cytoplasmic and nuclear proteins, energy metabolites) were set to 0.01$\times$ of their population-scale reference values in \texttt{phase3a\_spatial\_scaled.shy}. The extracellular species EPO and GCSF were not scaled, preserving their concentration in the solution bathing the cells; receptor species were likewise held at physiological values reflecting actual surface receptor numbers per cell. This choice is biologically motivated: cytokine concentrations are set by the bone marrow niche (an extensive extracellular space shared by many cells), while intracellular species scale with the volume of a single progenitor cell. At 0.01$\times$ initial conditions, the nuclear GATA1 place has an initial marking of approximately 8 tokens (molecules), consistent with published estimates of master transcription factor copy numbers in haematopoietic progenitor chromatin immunoprecipitation studies. The stochastic Gillespie semantics of Signal Hierarchical Petri Nets then simulate molecular events as discrete random occurrences, so that a single transcription event adds one molecule to the relevant place and produces a proportionally large relative fluctuation in the context of these small absolute numbers.

It is important to note that bimolecular transitions involving one internal species and one receptor species required careful treatment under this scaling: the transition rate was expressed in terms of the scaled internal species abundance while receptor occupancy was computed from the unscaled receptor concentration. This ensures that the effective concentration-driving rate of receptor-coupled signalling is physiologically consistent, rather than artificially reduced by applying the 0.01$\times$ factor to both sides.

\subsection{Simulation Ensemble}

Simulations were executed using the Gillespie stochastic algorithm as implemented in the Signal Hierarchical Petri Nets Viability Pane batch sweep interface. For the Phase B sweep, N=50 independent simulations were run for each of eight EPO concentrations (0.430, 0.440, 0.445, 0.449, 0.450, 0.451, 0.455, 0.460~mM) at fixed GCSF = 0.1~mM and fixed simulation duration $t = 3600$~s. GATA1 and PU.1 nuclear protein concentrations at $t = 3600$~s served as the primary commitment endpoints. Each simulation used a distinct random seed to ensure statistical independence. Simulation output included mean-over-replicates trajectories at 10001 equally spaced timepoints, per-replicate final states, and transition cumulative firing counts as functions of time.

Commitment fate was classified per replicate as erythroid if the final GATA1\_nuc exceeded PU.1\_nuc at $t = 3600$~s. This threshold-free classification is consistent with the phase portrait of the bistable attractor: replicates in the erythroid basin have GATA1\_nuc/PU.1\_nuc $>$ 1, while myeloid-committed replicates have GATA1\_nuc/PU.1\_nuc $<$ 1. Across all tested conditions, no replicate remained at equal levels at the final timepoint, consistent with the absence of a stable primed attractor at these copy numbers. The complete model specification and simulation configuration are available in the project repository.
